% !TeX root = ../教材题库_全部条目_转写工作稿.tex
% Source: XJ-B1-C04-TH-P098-Y1358
% 第4章 / 4.1 实数指数幂和幂函数 / 教材 p.98、99 / PDF 第105、106页
% OCR未校对：公式与文字均可能有误；未完成数学TeX转写；禁止直接发给学生
% 原题及相邻原解截图（核对用）：../题目原图/XJ-B1-C04-TH-P098-Y1358-01.webp, ../题目原图/XJ-B1-C04-TH-P098-Y1358-02.webp
% 本文件只在转写完成后作为题干使用；OCR 条目尚不是可靠题干。
\begin{bankitem}{p.98 思考（含紧邻说明）}{第4章 / 4.1 实数指数幂和幂函数 / 教材 p.98、99 / PDF 第105、106页}{XJ-B1-C04-TH-P098-Y1358}{OCR未校对：公式与文字均可能有误；未完成数学TeX转写；禁止直接发给学生}
\begin{ocrdraft}
[教材p.98 / PDF第105页]
\par
？思考我们知道，对于磊>0，当工是任意有理数时，酽都有了意义．
\par
在实际
\par
问题中，变量x可以是无理数，正数的无理数指数幂又怎么定义呢？
\par
下面以1萨为例来认识无理数指数幂．
\par
我们先来回顾确定无理数一1.414213“，的大小的过程：
\par
由于比1大，比2小；比1.4大，比1.5小；比L41大，比1.42小；比
\par
L414大，比1.415小；比L4142大，比1.4143小，
\par
“，如果需要，我们能够把
\par
[教材p.99 / PDF第106页]
\par
算到任意多位小数，即把它的大小范围估计到任意的精度，要多精确都可以达
\par
到，这样就算确定了它．
\par
当用一个小数近似代替后，利用有理数指数幂的基本不等式就可确定10万的
\par
取值范围．其过程如下表：
\par
的近似值范围
\par
1萨的近似值范围
\par
1.4<丿<1.5
\par
10L4<10、0<10L5
\par
1.41<《、/犭<1.42
\par
10L41<10万<10L12
\par
1.414<、伢<1.4巧
\par
101\ensuremath{\bullet}414<10、0<10
\par
I.415
\par
1.4142<《、/2<1.4143
\par
1()1\ensuremath{\bullet}4142<10、0<10L．
\par
借助计算器可得，10
\par
l.景142
\par
=25．95374，“>25，95，10
\par
l.景143
\par
=25．95971“，<25．96，
\par
于是25．95<10万<25．96，这样就获得了10、。的四位有效数字．
\par
由上可知，随着的近似值精度越高，1萨也会越接近一个实数，即10、。就是
\par
一个确定的实数．像这样，用磊的有理数次幂来逼近其无理数次幂，可以要多精确
\par
就有多精确．所以，任意正数a的无理数次幂就有了确定的意义．
\par
一般地，给定任意正数。对任意实数的次幂都有了定义．在幂的表
\par
达式中，a叫作底数，叫作指数．
\par
可以证明，有理数指数幂的前述运算法则，对实数指数幂仍然成立．
\par
即对于任意实数0、，均有：
\par
（3）伍b）厂=俨>0，b>0，r\ensuremath{\text{[OCR:EUR]}}R).
\par
类似地，我们有更一般的幂运算基本不等式：
\par
对任意的正实数和正实数。若>1，则">1；若磊<1，则aU<1.
\par
对任意的负实数和正实数山若>1，则"<1；若磊<1，则d>l.
\end{ocrdraft}
\end{bankitem}
